\title{FlashAttention-3: Fast and Accurate Attention with Asynchrony and Low-precision}

\begin{abstract}
As a foundational element in the Transformer architecture, attention remains a primary constraint for scaling large language models and supporting long-context tasks. \textsc{FlashAttention} introduced a strategy to accelerate attention computations on GPUs by reducing redundant memory access; nonetheless, its design does not leverage advanced functionalities in newer hardware, with \textsc{FlashAttention-2} utilizing merely 35\% of the available capacity on the H100 GPU. In this work, we present three key optimizations tailored for Hopper GPUs to further improve attention efficiency: utilizing the asynchronous operations of Tensor Cores and TMA to (1) enable concurrent computation and memory transfer through warp-specialization, (2) alternate block-level matrix multiplications and softmax calculations, and (3) implement block quantization and incoherent computation that exploits FP8 low-precision support in hardware. Our proposed method, \textsc{FlashAttention-3}, achieves 1.5–2.0$\times$ faster execution on H100 GPUs, attaining up to 740 TFLOPs/s (75\% of theoretical peak) using FP16, and approaching 1.2 PFLOPs/s with FP8 precision. Moreover, FP8 \textsc{FlashAttention-3} reduces numerical error by a factor of 2.6$\times$ compared to a conventional FP8-based attention mechanism.
\end{abstract}

\section{Introduction}
\label{sec:intro}

Within the Transformer architecture~\citep{vaswani2017attention}, the attention mechanism emerges as the central computational bottleneck, largely due to the fact that the calculation of self-attention scores between queries and keys incurs a quadratic cost relative to sequence length. Extending attention to accommodate longer context lengths promises to enable a variety of new functionalities, such as the modeling and reasoning across multiple extended documents~\citep{guo2021longt5,shaham2022scrolls,peng2023yarn}, or managing numerous files within large-scale codebases~\citep{roziere2023code, li2023starcoder}. Additionally, broader context allows the architecture to process new modalities, including high-resolution imagery~\citep{chen2022scaling}, audio~\citep{gulati2020conformer}, and video data~\citep{ho2022video}, as well as novel applications—ranging from handling extended user interaction histories~\citep{sun2019bert4rec} to orchestrating agent workflows spanning long time horizons~\citep{yao2022react}. As a result, there has been a surge in efforts to accelerate attention mechanisms for long-context scenarios, whether by introducing approximate methods~\citep{katharopoulos2020transformers,choromanski2020rethinking, tay2020efficient}, optimizing software implementations (\citep{rabe2021self,dao2022flashattention,kwon2023efficient}), or by exploring fundamentally different architectural approaches~\citep{peng2023rwkv,sun2023retentive,gu2023mamba}.

This study extends the foundational contributions of \citet{dao2022flashattention}, who advanced exact-attention algorithms by embedding insights from the GPU execution model and hardware properties into algorithmic design. Dao et al. presented \textsc{FlashAttention} in \citep{dao2022flashattention}, introducing an innovative tiling approach that accelerates attention computations by consolidating all attention steps into a single GPU kernel, thereby circumventing the need for intermediate accesses to slow global memory.

Subsequently, \citet{dao2023flashattention2} reorganized the methodology into \textsc{FlashAttention-2}, allowing parallelization across the sequence length axis and enabling the inner forward pass loop to process segments of the key and value matrices. This restructuring leads to better workload distribution and resource utilization on GPUs. Nonetheless, our analysis reveals that \textsc{FlashAttention-2} still falls short in terms of hardware utilization on current-generation GPUs when contrasted with highly tuned GEMM kernels; for instance, utilization rates are only about 35\%, compared to 80–90\% for GEMM on the Hopper H100 GPU. Part of this discrepancy can be explained by implementation factors, such as the continued reliance on Ampere-specific instructions rather than optimizing for Hopper architecture when leveraging Tensor Cores. Other works, including ThunkerKitten~\citep{spector2024thunder} and cuDNN 9~\citep{cudnn9}, demonstrate that adopting Hopper-optimized instructions alongside tile-based abstractions can significantly accelerate attention computation while also streamlining the software architecture.

At a deeper level, the algorithm underlying \textsc{FlashAttention-2} operates under a conventional synchronous paradigm and does not purposefully incorporate asynchronous operations or leverage low-precision arithmetic within its architecture. Asynchronous processing naturally emerges from the increasing presence of hardware components specifically optimized for machine learning tasks: distinct units, such as Tensor Cores for matrix multiplication and the Tensor Memory Accelerator (TMA) for memory transactions, function independently from general-purpose CUDA cores that handle logic, integer, and floating-point operations. The adoption of reduced-precision formats, like FP8 on Hopper and FP4 on Blackwell—building upon earlier transitions to FP16 (introduced in Pascal, 2017) and BF16 (debuted in Ampere, 2020)—has demonstrably allowed hardware to achieve two to four times greater throughput without increasing power consumption or silicon area. We detail how Hopper exploits these advancements in \cref{subsec:hardware}.
The primary technical obstacle, then, is to reengineer \textsc{FlashAttention-2} to effectively utilize these hardware optimizations: supporting asynchrony involves orchestrating overlapped execution of matmul and softmax phases despite their inherent data dependency, while enabling low-precision computation calls for sophisticated handling of quantization effects, particularly concerning rare feature activations present in LLMs~\citep{dettmers2208llm, sun2024massive}.

For this purpose, we introduce \textsc{FlashAttention-3}, which integrates and builds upon three novel concepts to enhance performance on modern GPU architectures:\footnote{While our results are contextualized using NVIDIA's Hopper architecture, the algorithm is applicable to any GPU platform equipped with strong asynchronous execution and support for low-precision arithmetic.}

\begin{enumerate}

\item \textbf{Producer-Consumer asynchrony:} We introduce a warp-level software pipelining method that leverages the decoupled execution capabilities of data transfer and Tensor Core computations by allocating producer and consumer roles to distinct warps. This strategy broadens the scope of latency hiding for memory accesses and instruction issuing within the algorithm.
\item \textbf{Hiding softmax under asynchronous block-wise GEMMs:} We conceal the relatively low-throughput, non-GEMM computations in softmax—such as floating-point fused multiply-add and exponentiation—by overlapping their execution with asynchronous block-wise GEMM operations using WGMMA instructions. To accomplish this, we modify the \textsc{FlashAttention-2} procedure to break select sequential dependencies between the softmax and GEMM phases. For instance, in our two-phase design, softmax processes one score matrix block at the same time that WGMMA asynchronously prepares the subsequent block.
\item \textbf{Hardware-accelerated low-precision GEMM:} We modify the forward computation routine to exploit FP8-capable Tensor Cores for GEMM, effectively achieving nearly twice the realized TFLOPs/s. This requires reconciliation of differing memory layout assumptions between WGMMA—specifically, the arrangement of FP32 accumulators and FP8 operands. To address the precision degradation inherent to FP8 computation, we employ block quantization and incoherent processing methods as compensatory measures.
\end{enumerate}

For empirical validation, we conduct benchmarks of \textsc{FlashAttention-3} on an H100 SXM5 GPU across varying parameter settings. The results indicate that (1) FP16 obtains a 1.5-2.0$\times$ acceleration over \textsc{FlashAttention-2} during the forward pass (achieving up to 740 TFLOPs/s), and a 1.5-1.75$\times$ speedup in the backward pass. (2) When operating in FP8, the method attains nearly 1.2 PFLOPs/s. (3) With longer sequence lengths, FP16 demonstrates superior performance, while FP8 remains competitive\footnote{Specifically, with a head dimension of 64, \textsc{FlashAttention-3} in FP8 exceeds, whereas for head dimensions 128 and 256 it matches in no causal masking scenarios and slightly trails with causal masking.} with NVIDIA's cuDNN-based leading attention implementation. Additionally, experiments confirm that \textsc{FlashAttention-3} in FP16 matches \textsc{FlashAttention-2} in terms of numerical precision, and surpasses conventional attention due to maintaining intermediate quantities (such as softmax rescaling) in FP32. Furthermore, \textsc{FlashAttention-3} in FP8 utilizing block quantization and incoherent computation is found to be 2.6$\times$ more accurate compared to conventional attention employing per-tensor quantization, particularly in instances with outlier features.

We are releasing \textsc{FlashAttention-3} under an open and permissive license\footnote{\textsc{FlashAttention-3}
  is available at \texttt{https://github.com/Dao-AILab/flash-attention}}, and intend to incorporate its support into the PyTorch and Hugging Face ecosystems to maximize its accessibility and usefulness for researchers and practitioners.

\section{Background: Multi-Head Attention and GPU Characteristics}
\label{sec:background}

\subsection{Multi-Head Attention}
\label{subsec:multi_head_attn}

Let $\mathbf{Q}, \mathbf{K}, \mathbf{V} \in \mathbb{R}^{N \times d}$ be the query, key and value input sequences associated to a single head, where $N$ is the sequence length and $d$ is the head dimension. Then the attention output $\mathbf{O}$ is computed as:

\begin{equation*}
  \mathbf{S} = \alpha \mathbf{Q} \mathbf{K}^\top \in \mathbb{R}^{N \times N}, \quad \mathbf{P} = \mathrm{softmax}(\mathbf{S}) \in \mathbb{R}^{N \times N}, \quad \mathbf{O} = \mathbf{P}\mathbf{V} \in \mathbb{R}^{N \times d},
\end{equation*}

Here, $\mathrm{softmax}$ is performed across each row, and it is common to choose the scaling parameter $\alpha = 1/\sqrt{d}$. To mitigate numerical issues associated with exponentiation, $\mathrm{rowmax}(\mathbf{S})$ is subtracted from $\mathbf{S}$. In the case of multi-head attention (MHA), distinct projection matrices for the query, key, and value are assigned to each head, enabling the computation to be efficiently parallelized across multiple attention heads and batches, ultimately yielding the complete output tensor.

Now let $\phi$ be a scalar loss function and let $\mathbf{d}(-) = \partial \phi / \partial (-)$ be notation for the gradient. Given the output gradient $\mathbf{dO} \in \mathbb{R}^{N \times d}$, we compute $\mathbf{dQ}$, $\mathbf{dK}$, and $\mathbf{dV}$ according to the chain rule as follows:

\begin{align*}
  \mathbf{dV} &= \mathbf{P}^\top \mathbf{dO} \in \mathbb{R}^{N \times d} \\
  \mathbf{dP} &= \mathbf{dO} \mathbf{V}^\top \in \mathbb{R}^{N \times N} \\
  \mathbf{dS} &= \mathrm{dsoftmax} (\mathbf{dP}) \in \mathbb{R}^{N \times N} \\
  \mathbf{dQ} &= \alpha \mathbf{dS} \mathbf{K} \in \mathbb{R}^{N \times d} \\
  \mathbf{dK} &= \alpha \mathbf{dS}^\top \mathbf{Q} \in \mathbb{R}^{N \times d},
\end{align*}

In this context, the relationship $\mathbf{d}s = (\mathrm{diag}(p) - p p^\top)\mathbf{d}p$ holds when expressing $p = \mathrm{softmax}(s)$, treating $p$ as a function of the vector $s$. The notation $\mathrm{dsoftmax}(\mathbf{dP})$ denotes the application of this rule independently to each row. As in the forward pass, the backward computation here can also be parallelized with respect to the number of heads and batches in MHA.

\subsection{GPU hardware characteristics and execution model}
\label{subsec:hardware}

We outline the elements of the GPU execution paradigm pertinent to \textsc{FlashAttention-3}, emphasizing the specifics of the NVIDIA Hopper architecture as an illustrative example of this framework.

\paragraph{Memory hierarchy:} The GPU architecture features a layered arrangement of memory resources, where higher storage capacities correspond to lower bandwidths (\cref{tab:gpu-hierarchy})\footnote{\citet{luo2024benchmarking} indicates that shared memory achieves a bandwidth of 128 bytes per clock cycle for each SM, which we scale by a factor of 132 SMs and the boost frequency of 1830 MHz.}. The global memory (GMEM), or HBM, is an off-chip DRAM resource available to every streaming multiprocessor (SM) within the device. Information stored in GMEM is automatically cached into an on-chip L2 cache to facilitate quicker access. Each SM is further equipped with its own relatively small, on-chip, high-bandwidth, and heavily banked shared memory (SMEM), which is under explicit software control. Finally, the register file resides within every SM, providing the fastest and most localized storage.

\paragraph{Thread hierarchy:} The structure of GPU programming relies on a hierarchical organization of its computational elements, termed threads. Arranged from the most granular to the most aggregate level, this hierarchy includes individual threads, warps consisting of 32 threads each, warpgroups formed by grouping 4 consecutive warps, threadblocks (also referred to as cooperative thread arrays or CTAs), threadblock clusters as introduced in Hopper architecture, and ultimately, grids at the highest level.

The relationship between these two hierarchies is tightly coupled. Threads belonging to a single CTA are concurrently executed on the same SM, and CTAs grouped within one cluster are run together on the same GPC. While SMEM is accessible to every thread inside a CTA, each individual thread is limited to a maximum of 256 private registers (RMEM).

\paragraph{Asynchrony and warp-specialization:}

GPUs achieve high throughput by leveraging parallelism and asynchronous operations, which serve to mask memory and computation delays. In the Hopper architecture, asynchronous memory transfers between GMEM and SMEM are managed by a specialized hardware component known as the Tensor Memory Accelerator (TMA) \cite[\S7.29]{cuda}. Additionally, Hopper distinguishes itself from earlier generations like Ampere by equipping its Tensor Core—accessible through the warpgroup-wide WGMMA instruction \cite[\S9.7.14]{ptx}—with asynchronous capabilities, enabling it to directly fetch data from shared memory.

The presence of hardware-level asynchrony enables the use of warp-specialized kernels, wherein the warps within a CTA are segregated into either producer or consumer categories—dedicating themselves solely to either data transfer operations or computational tasks. Such an approach generally enhances the compiler’s capacity to produce superior instruction schedules \citep{warp-specialization-2011}. Moreover, Hopper facilitates flexible register allocation among warpgroups through the use of \verb|setmaxnreg| \cite[\S9.7.17.1]{ptx}, permitting warps engaged in MMAs to access a greater portion of RMEM compared to those limited to TMA operations, which require just a single active thread.

\paragraph{Low-precision number formats:}
\label{sec:low-precision-gpu}
Contemporary GPUs incorporate dedicated hardware specifically designed to expedite computations using reduced-precision formats. As an illustration, the WGMMA instruction leverages the FP8 Tensor Cores available on Hopper architecture, resulting in twice the per-SM throughput relative to computations performed in FP16 or BF16.

Successfully utilizing FP8 WGMMA requires careful consideration of the operand layout restrictions. Consider the matrix multiplication $A \times B^{\top}$, where $A$ is of dimension $M\times K$ and $B$ is $N\times K$. An operand is described as \emph{mn-major} if the data is arranged contiguously along the outer $M$ or $N$ dimension, whereas it is termed \emph{k-major} if contiguity is along the inner $K$ dimension. In the context of FP16 WGMMA, both mn-major and k-major layouts are valid for input operands stored in SMEM. In contrast, FP8 WGMMA restricts support to only the k-major operand layout. Additionally, in applications like attention—where chaining multiple GEMMs within a single kernel is desired—the conflicting layouts between the FP32 accumulator and FP8 operands impede the seamless invocation of successive FP8 WGMMAs.

Regarding attention mechanisms, these layout constraints necessitate specific adjustments to how an FP8 algorithm is constructed, as detailed in \cref{sec:algofp8}.

\subsection{Standard Attention and Flash Attention} Consistent with \citet{dao2022flashattention}, we use the term \textbf{standard attention} to refer to a GPU-based attention computation where the intermediate matrices $\mathbf{S}$ and $\mathbf{P}$ are explicitly stored in HBM. The key innovation introduced by \textsc{FlashAttention} was to apply a localized version of the softmax reduction, thereby mitigating the costly memory accesses involved in storing and loading these intermediates, and to integrate the attention mechanism within a single GPU kernel. This local softmax process encompasses lines \ref{code-ws:softmax_start}-\ref{code-ws:softmax_end} of the consumer mainloop as specified in \cref{alg:flash3_wgmma_ws_only}, along with the rescaling operations applied to blocks of $\mathbf{O}$. A straightforward derivation confirming that this approach correctly produces $\mathbf{O}$ is provided in \cite[\S 2.3.1]{dao2023flashattention2}.

\section{FlashAttention-3: Algorithm}
\label{sec:algo}

This section outlines the \textsc{FlashAttention-3} algorithm. For clarity, we limit our discussion to the forward pass here, while the method for the backward pass is presented in~\cref{sec:algo_ws_bwd}. We begin by detailing how warp-specialization combined with a circular SMEM buffer can be incorporated into the foundational framework of \textsc{FlashAttention-2}. Next, we illustrate how leveraging the asynchrony inherent in WGMMA enables the creation of a two-stage pipeline that overlaps GEMM and softmax operations. Lastly, we discuss the adjustments necessary for FP8 support, addressing both layout requirements and improvements in accuracy achieved through block quantization and processing without coherence.

\subsection{Producer-Consumer asynchrony through warp-specialization and pingpong scheduling}
\label{sec:algo_ws}

\paragraph{Warp-specialization}
Similar to \textsc{FlashAttention-2}, the forward propagation in \textsc{FlashAttention-3} is highly parallelizable across batch size, number of attention heads, and length of the query sequence. Consequently, the algorithm can be effectively described at the CTA level, where it processes a specific tile $\mathbf{Q}_i$ from the query matrix in order to generate the associated output tile $\mathbf{O}_i$. For clarity, we initially present the warp-specialization approach utilizing a circular SMEM buffer, excluding any overlap between GEMM and softmax computations. Here, let $d$ denote the head dimension and $N$ the length of the sequence. By selecting a query block size $B_r$, we partition $\mathbf{Q}$ into $T_r = \lceil \frac{N}{B_r} \rceil$ blocks labeled as $\mathbf{Q}_1, .., \mathbf{Q}_{T_r}$.

\begin{algorithm}[H]
    \caption{\small\label{alg:flash3_wgmma_ws_only}\textsc{FlashAttention-3} forward pass \textbf{without} intra-consumer overlapping -- CTA view}
    \begin{algorithmic}[1]
\REQUIRE Matrices $\mathbf{Q}_i \in \mathbb{R}^{B_r \times d}$, $\mathbf{K}, \mathbf{V} \in \mathbb{R}^{N \times d}$ stored in HBM, key block size $B_c$, with the number of blocks $T_c = \lceil \frac{N}{B_c} \rceil$.
\STATE Set up pipeline management object to coordinate barrier synchronization, using an $s$-stage circular shared memory buffer.
\IF {current warpgroup is a producer}
\STATE Release a predefined quantity of registers.
\STATE Fetch $\mathbf{Q}_i$ from HBM into shared memory.
\STATE Once completed, signal to the consumer that $\mathbf{Q}_i$ is ready.
\FOR{$0 \le j < T_c$}
    \STATE Await consumption of the $(j\,\%\,s)$th buffer stage.
    \STATE Fetch $\mathbf{K}_j, \mathbf{V}_j$ from HBM into the $(j\,\%\,s)$th slot of shared memory buffer.
    \STATE After loading, signal to consumers that $\mathbf{K}_j, \mathbf{V}_j$ are available.
\ENDFOR
\ELSE 
\STATE Allocate an appropriate number of registers based on the number of consumer warps.
\STATE Initialize on-chip: $\mathbf{O}_i = (0) \in \mathbb{R}^{B_r \times d}$, $\ell_i, m_i = (0), (-\infty) \in \mathbb{R}^{B_r}$.
\STATE Wait until $\mathbf{Q}_i$ is accessible in shared memory.
\FOR{$0 \le j < T_c$}
\STATE Wait until $\mathbf{K}_j$ is ready in shared memory.
\STATE Perform matrix multiplication $\mathbf{S}_i^{(j)} = \mathbf{Q}_i \mathbf{K}_j^T$ (SS-GEMM), then synchronize. \label{alg:ws_only_gemm1}
\STATE Save current $m_i$ as $m_i^{\mathrm{old}}$, update $m_i = \mathrm{max}(m_i^{\mathrm{old}}, \mathrm{rowmax}(\mathbf{S}_i^{(j)}))$. \label{code-ws:softmax_start}
\STATE Compute $\widetilde{\mathbf{P}}_i^{(j)} = \mathrm{exp}(\mathbf{S}_i^{(j)} - m_i)$; update $\ell_i = \mathrm{exp}(m_i^{\mathrm{old}} - m_i) \ell_i + \mathrm{rowsum}(\widetilde{\mathbf{P}}_i^{(j)})$. \label{code-ws:softmax_end}
\STATE Wait for $\mathbf{V}_j$ to finish loading to shared memory.
\STATE Update $\mathbf{O}_i = \mathrm{diag}(\mathrm{exp}(m_i^{\mathrm{old}} - m_i))^{-1} \mathbf{O}_i + \widetilde{\mathbf{P}}_i^{(j)} \mathbf{V}_j$ (RS-GEMM), then synchronize. \label{alg:ws_only_gemm2}
\STATE Free the $(j\,\%\,s)$th buffer stage for subsequent producer operations.
\ENDFOR
\STATE Normalize the output by computing $\mathbf{O}_i = \mathrm{diag}(\ell_i)^{-1} \mathbf{O}_i$ and calculate $L_i = m_i + \log(\ell_i)$.
\STATE Store both $\mathbf{O}_i$ and $L_i$ back to HBM in the $i$th position of $\mathbf{O}$ and $L$.
\ENDIF
\end{algorithmic}
\end{algorithm}

In our realization of \cref{alg:flash3_wgmma_ws_only} on the Hopper architecture, we employ \verb|setmaxnreg| for register (de)allocation, utilize TMA to load $\mathbf{Q}_i$ and $\{ \mathbf{K}_j, \mathbf{V}_j \}_{0 \leq j < T_c}$, and rely on WGMMA for executing GEMMs within the main loop of the consumer, designating the SS or RS prefix to specify whether the initial operand is fetched from shared memory or the register file, respectively. When analyzing the control flow of \cref{alg:flash3_wgmma_ws_only}, it is important to recognize that TMA loads are asynchronous and therefore, launching these loads does not depend on the completion of previous ones. Additionally, within the producer's main loop, no explicit waits occur during the initial $s$ iterations, as these cycles are dedicated to populating the buffer.

\paragraph{Pingpong scheduling} The asynchronous execution characteristics of both WGMMA and TMA, combined with warp specialization, enable overlapping the softmax calculation of one warpgroup concurrently with the GEMM execution of another. This is motivated by the observation that, on contemporary hardware accelerators, operations aside from matmul are processed at substantially lower throughputs. For instance, on the H100 SXM5 GPU, FP16 matmul achieves a peak of 989 TFLOPS, whereas special functions such as the exponential achieve only 3.9 TFLOPS\footnote{According to the CUDA programming guide, each streaming multiprocessor (SM) can complete 16 special function operations per clock. This is calculated as 16 operations $\times$ 132 SMs $\times$ 1830 MHz, yielding 3.9 TFLOPS for special functions.}, which are crucial for implementing softmax. In the attention forward pass, when operating in FP16 and using a head dimension of 128, the ratio of matmul to exponential FLOPS is 512:1, but the hardware throughput for exponential is 256 times less, making the total time taken for exponentials reach up to 50\% of that for matmul. This disparity grows with FP8 computation: matmul throughput is doubled, but the performance of special functions, such as exponentials, does not increase.

Because the exponential function is executed by a dedicated hardware component—the multi-function unit—it is preferable for this computation to occur concurrently with the Tensor Cores' matmul operations. To achieve this overlap, we introduce synchronization barriers (\texttt{bar.sync} instructions), ensuring that the GEMM operations (GEMM1, corresponding to $\mathbf{P} \mathbf{V}$ in the current iteration, and GEMM0, corresponding to $\mathbf{Q} \mathbf{K}^\top$ for the following iteration) in warpgroup 1 are scheduled before those in warpgroup 2. Consequently, while warpgroup 1 undertakes the softmax computation, warpgroup 2 proceeds with its GEMMs. The process then alternates: warpgroup 2 computes softmax as warpgroup 1 returns to GEMMs. This technique, referred to as ``pingpong'' scheduling, is depicted in~\cref{fig:pingpong_scheduling}. 

Although the practical implementation of pingpong scheduling does not always perfectly align with the schematic shown, empirical results indicate performance benefits—such as an increase in throughput from 570 TFLOPS to 620–640 TFLOPS for FP16 forward pass with a head dimension of 128 and a sequence length of 8192.

\paragraph{Attention variants} In the case of multi-query attention~\citep{shazeer2019fast} and grouped query attention~\citep{ainslie2023gqa}, we adopt the methodology from \textsc{FlashAttention-2}, modifying tensor indices to prevent redundant copies of $\mathbf{K}$ and $\mathbf{V}$ within HBM.

\subsection{Intra-warpgroup overlapping GEMMs and softmax}

It is possible to achieve partial instruction overlap between softmax operations and GEMM computations within a single warpgroup. In the following, we outline a specific method that facilitates this overlap.

In the attention algorithm, operations within the inner loop (main loop) have sequential dependencies that impede parallelization within a single iteration. For example, (local) softmax (lines \ref{code-ws:softmax_start} to \ref{code-ws:softmax_end}) relies on the output $\mathbf{S}_i^{(j)}$ of the first GEMM, while the second GEMM takes its result $\widetilde{\mathbf{P}}_i^{(j)}$ as an operand. Indeed, the wait statements in lines \ref{alg:ws_only_gemm1} and \ref{alg:ws_only_gemm2} of \cref{alg:flash3_wgmma_ws_only} serialize the execution of softmax and GEMMs. However, we can break these dependencies by pipelining across iterations through additional buffers in registers. Pursuing this idea, we propose the following two-stage\footnote{Note that the number of stages of the overlapping scheme is bounded by, but need not equal, the number $s$ of stages in the circular SMEM buffer.} GEMM-softmax pipelining algorithm:

\begin{algorithm}[H]
    \caption{\small\label{alg:flash3_wgmma}\textsc{FlashAttention-3} consumer warpgroup forward pass}
    \begin{algorithmic}[1]
      \REQUIRE  Given matrices $\mathbf{Q}_i \in \mathbb{R}^{B_r \times d}$ and $\mathbf{K}, \mathbf{V} \in \mathbb{R}^{N \times d}$ stored in HBM, along with the key block size $B_c$ and $T_c = \lceil \frac{N}{B_c} \rceil$.
      \STATE Adjust register allocation according to the specified number of consumer warps.
      \STATE Initialize the on-chip matrices and vectors: set $\mathbf{O}_i = (0) \in \mathbb{R}^{B_r \times d}$, $\ell_i = (0) \in \mathbb{R}^{B_r}$, and $m_i = (-\infty) \in \mathbb{R}^{B_r}$.
      \STATE Await the transfer of $\mathbf{Q}_i$ and the first key block $\mathbf{K}_0$ into shared memory.
      \STATE Use WGMMA to compute $\mathbf{S}_{\mathrm{cur}} = \mathbf{Q}_i \mathbf{K}_0^T$; commit the operation and ensure completion before proceeding.
      \STATE Free up the initial stage ($0$th) in the $\mathbf{K}$ buffer.
      \STATE Based on $\mathbf{S}_{\mathrm{cur}}$, determine $m_{i}$, $\tilde{\mathbf{P}}_{\mathrm{cur}}$, and $\ell_{i}$; apply scaling to $\mathbf{O}_i$ accordingly.
      \FOR{$1 \le j < T_c - 1$}
        \STATE Wait until block $\mathbf{K}_j$ is present in shared memory.
        \label{code:mainloop_start}
        \STATE Use WGMMA to calculate $\mathbf{S}_{\mathrm{next}} = \mathbf{Q}_i \mathbf{K}_j^T$; initiate the computation without waiting for its completion. \label{code:first_wgmma}
        \STATE Wait for the previous value block $\mathbf{V}_{j-1}$ to be loaded into shared memory.
        \STATE With WGMMA, compute the update $\mathbf{O}_{i} = \mathbf{O}_{i} + \tilde{\mathbf{P}}_{\mathrm{cur}} \mathbf{V}_{j-1}$; start the operation but proceed without waiting. \label{code:second_wgmma}
        \STATE Ensure the $\mathbf{Q}_i \mathbf{K}_j^T$ operation with WGMMA has completed.
        \STATE Derive $m_{i}$, $\tilde{\mathbf{P}}_{\mathrm{next}}$, and $\ell_{i}$ using $\mathbf{S}_{\mathrm{next}}$ as input. \label{code:softmax}
        \STATE Wait for completion of the $\tilde{\mathbf{P}}_{\mathrm{cur}} \mathbf{V}_{j-1}$ computation; subsequently, rescale $\mathbf{O}_i$.
        \STATE Release the $(j\,\%\,s)$th buffer segment for $\mathbf{K}$ and the $(j-1\,\%\,s)$th for $\mathbf{V}$.
        \STATE Move $\mathbf{S}_{\mathrm{next}}$ data into $\mathbf{S}_{\mathrm{cur}}$.
        \label{code:mainloop_end}
      \ENDFOR \STATE Wait for the final value block $\mathbf{V}_{T_c - 1}$ to be available in shared memory.
      \STATE Update the output matrix:
      $\mathbf{O}_{i} = \mathbf{O}_{i} + \tilde{\mathbf{P}}_{\mathrm{last}} \mathbf{V}_{T_c - 1}$ via WGMMA; ensure the operation completes before exiting.

\STATE Epilogue: Adjust $\mathbf{O}_{i}$ using $m_{i}$. Derive $L_{i}$ utilizing $m_{i}$ and $\ell_{i}$.
      Store $\mathbf{O}_{i}$ and $L_{i}$ in HBM as the $i$-th segments of $\mathbf{O}$ and $L$.
    \end{algorithmic}
  \end{algorithm}

\cref{alg:flash3_wgmma} substitutes for the consumer section in \cref{alg:flash3_wgmma_ws_only} to assemble the full \textsc{FlashAttention-3} method targeting FP16 computations. Broadly, WGMMA is utilized as shorthand to refer to asynchronous GEMM execution. Inside the mainloop (see lines \ref{code:mainloop_start} through \ref{code:mainloop_end}), the second WGMMA invocation in iteration $j$ (line \ref{code:second_wgmma}) is executed concurrently with the softmax computations for the subsequent iteration $j+1$ (line \ref{code:softmax}).

Although the pipelined design depicted above promises theoretical enhancements in throughput, several pragmatic challenges must be addressed:
\paragraph{Compiler reordering} While the pseudocode outlines an idealized sequence of execution, the actual ordering of instructions is often modified by the NVCC compiler as part of its optimization strategies. Such reordering can interfere with the intended overlap between WGMMA and non-WGMMA operations, leading to potential loss of intended pipelining efficiency or unpredictable effects. Inspection of the compiled SASS code reveals that, in practice, instruction overlap is generally realized as designed (Section~\ref{sec:2-stage-sass}).
\paragraph{Register pressure} Achieving high efficiency requires careful management to avoid register spilling. The 2-stage pipeline inherently needs more registers for holding intermediate values and for coordinating transitions between pipeline stages. In particular, it necessitates that the variable $\mathbf{S}_{\mathrm{next}}$ is allocated in registers, adding an overhead of $B_r \times B_c \times \text{sizeof}(\text{float})$ registers for each threadblock. This surge in register utilization can compete directly with optimizations like increasing thread block size, which itself demands more registers. Therefore, a balanced approach, informed by thorough profiling, is required to manage these competing demands.
\paragraph{3-stage pipelining} Building upon the previously described 2-stage approach, a 3-stage variant can be constructed to introduce overlap between a second WGMMA operation and the softmax computation. While this expansion holds promise for improved utilization of Tensor Cores, it brings about even higher register requirements due to the extra stage. This heightens the challenge of reconciling the trade-offs between pipeline depth and tile dimensions. The precise details and empirical findings regarding the 3-stage algorithm are documented in ~\cref{sec:3-stage}.

\subsection{Low-precision with FP8}
\label{sec:algofp8}

\textbf{Efficiency: layout transformations.} Executing the forward pass of \textsc{FlashAttention-3} using FP8 precision introduces unique difficulties related to layout conformity that are not present when operating in FP16.

To begin with, it is important to observe that the input tensors $\mathbf{Q}$, $\mathbf{K}$, and $\mathbf{V}$ are commonly stored such that the head dimension is contiguous in memory. However, to comply with the k-major layout requirement of FP8 WGMMA for the second GEMM, it is necessary that $\mathbf{V}$—specifically, the tiles of $\mathbf{V}$ loaded into SMEM—be contiguous along the sequence length dimension. Because the TMA load operation preserves the original contiguous axis, we face a choice: either (1) perform a transposition of $\mathbf{V}$ in GMEM as a preprocessing action, or (2) apply an in-kernel transpose to the $\mathbf{V}$ tiles once they are loaded into SMEM. For the first approach, there are two possible implementations: (1a) the transpose could be fused into the epilogue of the previous step (for example, after rotary embedding), or (1b) a separate transpose kernel could be invoked before processing\footnote{A highly optimized transpose kernel can operate close to the hardware bandwidth limit~\citep{colfax_cutlass_transpose_2024}.}, thus swapping the strides between the sequence length and head dimensions. Nonetheless, approach (1a) poses significant challenges for integration into a generic library, whereas (1b) results in considerable inefficiency, especially in inference settings dominated by memory bandwidth.

For FP8 \textsc{FlashAttention-3}, we choose option (2). During the in-kernel transpose, we exploit the LDSM (\verb|ldmatrix|) and STSM (\verb|stmatrix|) instructions. These operations enable a warp of threads to collaboratively move data from SMEM to RMEM and vice versa, processing memory in 128-byte increments.\footnote{While the PTX documentation describes LDSM/STSM as handling the transfer of $8 \times 8$ matrices containing 16-bit values \cite[\S 9.7.13.4.15-16]{ptx}, it is possible to combine pairs of 8-bit values for use with LDSM/STSM in FP8 settings. Nevertheless, the transpose variants of LDSM/STSM are not able to split apart these packed 8-bit values, requiring certain register shuffles between LDSM and STSM to accomplish the actual tile-level transpose; further specifics are omitted.} LDSM and STSM instructions are both register-friendly, which permits execution within the producer warpgroup, and can reconfigure layouts during memory copying for transposition. In addition, after the initial iteration, it is possible to schedule the transpose for the next $\mathbf{V}$ tile concurrently with the two WGMMAs that correspond to the previous $\mathbf{V}$ and the current $\mathbf{K}$ tile.

Second, our analysis reveals that, in contrast to FP16, the register-resident FP32 accumulator resulting from an FP8 WGMMA has a distinct memory layout compared to operand A. This distinction is illustrated in \cref{fig:rmem_accum} and \cref{fig:rmem_operand}, which display partial layouts with the register allocation per thread given in sequential order. Through the application of byte permutation instructions, we are able to convert the accumulator generated by the initial WGMMA into a representation that aligns with the input format for a subsequent WGMMA. This transformed layout also conforms to the arrangement of the $\mathbf{V}$ tile following the in-kernel transpose. More concretely, referring to \cref{fig:rmem_accum}, the required permutation reorganizes the registers to follow the pattern
$$\{ \verb|d0 d1 d4 d5 d2 d3 d6 d7| \},$$
and this reordering is uniformly applied to every group of 8 bytes. On a structural level, this procedure effectively permutes the columns within the $\mathbf{P}$ tile (for instance, shifting columns $0189$ to the leading positions). To ensure WGMMA produces the appropriate output tile, the in-kernel transpose can be directed to write the $\mathbf{V}$ tile according to a corresponding row permutation.\footnote{This added flexibility enabled by executing the transpose within the kernel removes the need for shuffle operations to redistribute register values between threads—a step that was necessary in our previous work~\citep{colfax_fp8_flashattention_2024}.}

\textbf{Accuracy: block quantization and incoherent processing.} The FP8 (e4m3) data type allocates 3 bits to the mantissa and 4 bits to the exponent, which leads to greater numerical imprecision in comparison to FP16 or BF16. In large-scale models, it is common to encounter a few outlier values that substantially exceed the rest in magnitude~\citep{dettmers2208llm,
  sun2024massive}, complicating the quantization process. To address this, per-tensor scaling is generally adopted~\citep{micikevicius2022fp8}, wherein a single scaling factor is maintained for each tensor (for example, one each for $\mathbf{Q}$, $\mathbf{K}$, and $\mathbf{V}$).
To address the increased numerical error associated with FP8 in attention computations, we introduce two main strategies:

\begin{enumerate}

\item \textbf{Block quantization}: In this approach, a single scalar value is stored per block, whereby for each of $\mathbf{Q}$, $\mathbf{K}$, and $\mathbf{V}$, the tensors are partitioned into segments of dimensions $B_r \times d$ or $B_c \times d$, and quantization is conducted individually on each segment. This quantization process can be integrated with preceding operations in the attention computation—such as rotary embedding—without causing additional latency, since the rotary embedding step is constrained by memory bandwidth. Given that the \textsc{FlashAttention-3} algorithm processes data in blocks by design, it is feasible to appropriately rescale every block of $\mathbf{S}$ to accommodate for block-based quantization with no extra computation required.

\item \textbf{Incoherent processing}: To mitigate the impact of outliers, we pre-multiply $\mathbf{Q}$ and $\mathbf{K}$ by a randomly sampled orthogonal matrix $\mathbf{M}$ prior to applying FP8 quantization. Because $\mathbf{M}$ is orthogonal, $\mathbf{M} \mathbf{M}^\top = I$, ensuring that $(\mathbf{Q} \mathbf{M})(\mathbf{K} \mathbf{M})^\top = \mathbf{Q} \mathbf{K}^\top$; therefore, the multiplication leaves the output of the attention mechanism unchanged. The effect of this transformation is to distribute outlier values more uniformly, since each entry in $\mathbf{Q} \mathbf{M}$ or $\mathbf{K} \mathbf{M}$ becomes a stochastic combination of the original values in $\mathbf{Q}$ or $\mathbf{K}$, thus helping decrease quantization error. In our implementation, following the recommendations of \citet{chee2024quip} and~\citet{tseng2024quip}, $\mathbf{M}$ is taken as the product of random diagonal matrices whose entries are $\pm 1$ and a Hadamard matrix. Such a structure permits multiplication to be performed in $O(d \log d)$, improving efficiency over the $O(d^2)$ general case, and, furthermore, this transformation can be integrated with rotary embedding operations without incurring additional computational expense.
\end{enumerate}

We empirically demonstrate in \cref{sec:numerical_error} that these two approaches decrease numerical errors by a factor of up to 2.6$\times$.

\section{Empirical Validation}
\label{sec:exp}

Primitives provided by CUTLASS~\citep{Thakkar_CUTLASS_2023}, specifically WGMMA and TMA abstractions, are utilized to construct \textsc{FlashAttention-3} and to assess its performance in terms of efficiency and accuracy.

\begin{itemize}

\item \textbf{Attention benchmarking.} We evaluate the performance of \textsc{FlashAttention-3} at various sequence lengths, contrasting its execution time against the standard PyTorch implementation, \textsc{FlashAttention-2}, as well as a Triton variant of \textsc{FlashAttention-2} that utilizes H100-specific instructions, and a vendor-supplied cuDNN-optimized implementation for H100 GPUs. Our findings demonstrate that \textsc{FlashAttention-3} outperforms \textsc{FlashAttention-2} by a factor of up to 2.0$\times$, and surpasses the Triton-optimized \textsc{FlashAttention-2} by 1.5$\times$. Additionally, \textsc{FlashAttention-3} achieves throughput as high as 740 TFLOPs/s, equating to 75\% of the H100 GPUs’ peak theoretical throughput.

\item \textbf{Ablation experiments.} Through controlled studies, we show that enhancements such as warp-specialization and GEMM-softmax pipelining are central to the observed acceleration in \textsc{FlashAttention-3}.

\item \textbf{FP8 attention accuracy.} Our analysis indicates that employing block quantization and incoherent processing decreases the numerical errors in FP8 \textsc{FlashAttention-3} by 2.6$\times$.

\subsection{Benchmarking Attention}
\label{subsec:benchmark_attn}

We evaluate the performance of various attention mechanisms by recording their execution time on an H100 80GB SXM5 GPU under different scenarios (presence or absence of causal masking, head dimensions set to 64 or 128), utilizing FP16 input representations. Performance data is summarized in~\cref{fig:benchmark_attn_fwd} and~\cref{fig:benchmark_attn_bwd}. The findings demonstrate that \textsc{FlashAttention-3} outperforms \textsc{FlashAttention-2} by approximately 1.5–2.0$\times$ in forward computations and by 1.5–1.75$\times$ in backward computations. Relative to the conventional attention mechanism, \textsc{FlashAttention-3} achieves speedups ranging from 3$\times$ to 16$\times$. For sequence lengths of 1k or higher, \textsc{FlashAttention-3} not only matches but often exceeds the speed of the specialized cuDNN library, a proprietary tool optimized for the H100 GPU architecture.

\paragraph{Benchmark settings:} We vary the sequence length as 512, 1k, ..., 16k, and set batch size so that the total number of tokens is 16k. We set the hidden dimension to 2048, and head dimension to be either 64, 128, or 256 (i.e., 32 heads, 16 heads, or 8 heads). To calculate the FLOPs of the forward pass, we use:

\begin{equation*}
  4 \cdot \text{seqlen}^2 \cdot \text{head dimension} \cdot \text{number of heads}.
\end{equation*}

Applying causal masking requires halving this value, as roughly 50% of the computations are performed. To estimate the FLOPs required for the backward pass, we scale the forward pass FLOPs by a factor of 2.5, reflecting the fact that the forward operation involves 2 matrix multiplications, while the backward operation involves 5—this difference arises from recomputation.

Additionally, we assess the forward pass runtime of FP8 in comparable conditions. Results corresponding to a headdim of 256 are depicted in~\cref{fig:benchmark_attn_fp8}, while comprehensive results can be found in \cref{sec:benchmark_attn_fp8_full}.

\subsection{Ablation Study: 2-Stage Pipelining Experiments}
\label{sec:2-stage-pipelining-experiments}

We conduct an ablation study on the 2-stage WGMMA-softmax pipeline and on warp-specialization within the context of non-causal FP16 \textsc{FlashAttention-3}, using the fixed configuration $\{\text{batch}, \text{seqlen}, \text{nheads}, \text{hdim}\} = \{ 4, 8448, 16, 128\}$. As indicated in~\cref{table:ablation_pipelining}, the introduction of algorithmic enhancements—specifically, leveraging asynchrony via warp-specialization and enabling the overlap of GEMM and softmax computation—delivers a substantial performance boost, increasing throughput from 570 to 661 TFLOPs.

\subsection{Numerical Error Validation}
\label{sec:numerical_error}

Given recent attention to the numerical error of \textsc{FlashAttention}~\citep{golden2024flash}, we evaluate and contrast the numerical precision of \textsc{FlashAttention-2}, \textsc{FlashAttention-3}, and a conventional attention implementation by benchmarking them against a reference in FP64. In order to mimic the presence of outlier features and activations as observed in LLMs~\citep{dettmers2208llm, sun2024massive}, we sample the elements of $\mathbf{Q}, \mathbf{K}, \mathbf{V}$ according to the following distribution:

\begin{equation*}
  \mathcal{N}(0, 1) + \mathcal{N}(0, 100) \cdot \mathrm{Bernoulli}(0.001).
\end{equation*}

Each element follows a normal distribution with mean zero and unit standard deviation, except that 0.1\% of the entries include an additional independent Gaussian component with a standard deviation of 10. Subsequently, we compute the root mean squared error (RMSE) as shown in~\cref{table:numerical_error}. When using FP16, both \textsc{FlashAttention-2} and \textsc{FlashAttention-3} demonstrate a 1.7$\times$ reduction in RMSE relative to the default approach, attributed to storing intermediate softmax computations in FP32. For the baseline FP8 attention, a per-tensor scaling strategy is adopted, with the matmul accumulator in FP32 while the intermediate softmax values remain in FP16. Owing to block quantization and disjoint (incoherent) computation, \textsc{FlashAttention-3} implemented in FP8 achieves a RMSE that is 2.6$\times$ lower than this baseline.

\section{Dicussion, Limitations, Conclusion}
\label{sec:discussion}

In this work, we show that leveraging novel programming strategies alongside advancements in hardware—specifically asynchrony and low-precision operations—significantly enhances both the performance and accuracy of attention mechanisms using \textsc{FlashAttention-3}. Our approach achieves attention speedups of 1.5–2.0$\times$ over \textsc{FlashAttention-2}, while decreasing FP8 numerical errors by a factor of 2.6 in comparison to typical per-tensor quantization methods. Several aspects remain to be improved upon, including: further optimizations for LLM inference, adopting a persistent kernel approach for the FP8 kernel,\footnote{In the context of our benchmarks, the FP16 version of \textsc{FlashAttention-3} incorporates a persistent kernel and a load balancing mechanism, whereas these features are not present in the FP8 variant. This partly accounts for FP8 \textsc{FlashAttention-3}'s relatively lower performance with shorter sequence lengths and causal masking as compared to FP8 cuDNN kernels.} and more comprehensive investigation into the role of low-precision attention during large-scale training regimes. While our current evaluation centers on Hopper GPUs, we anticipate that the principles and methods presented here are broadly transferable to a range of other hardware accelerators. Ultimately, we believe the development of a faster, more accurate attention primitive could spur progress and foster innovation in tasks demanding extended context.

\end{document}